\subsection{Lean 4 Formalisation}
The entire mathematical model defined in Section 4 has been mechanised in Lean 4. The project compiles without any use of the \texttt{sorry} or \texttt{admit} keywords, ensuring full rigorous verification.

\subsection{Code Snippets}
The transition relation is defined as:
\begin{lstlisting}[language=Lean]
def Transition {V Q A B Gamma : Type} (scen : Scenario V Q A B Gamma) (assumptions : Gamma) (s₁ s₂ : GlobalState V Q A B) : Prop :=
  (s₁.zone, s₂.zone) ∈ scen.edges ∧
  scen.guard (s₁.zone, s₂.zone) s₁.cap s₁.caregiver s₁.barrier assumptions ∧
  s₂.cap = s₁.cap ∧
  s₂.caregiver = s₁.caregiver ∧
  s₂.barrier = s₁.barrier
\end{lstlisting}

The core safety theorem is stated as:
\begin{lstlisting}[language=Lean]
theorem reachability_invariant_equivalence {V Q A B Gamma : Type} (scen : Scenario V Q A B Gamma) (assumptions : Gamma) :
    Safe scen assumptions ↔ ¬InvariantDeficit scen assumptions
\end{lstlisting}

\subsection{Mechanisation Statistics}
The formalisation comprises approximately 10 Lean modules, formalising finite transition systems, reachability, safety invariants, interventions, graph cuts, robustness, and conservative abstraction. The codebase includes the running case study of the two-storey house.
